
% %%%%%%%%%%%%%%%%%%%%%%%%%%%%%%%%%%%%%%%%%%%%%%%%%%%%%%%%%%%%%%%%%%%%%%%%
% %%%%%%%%%%%%%%%%%%%%%%%%%%%%%%%%%%%%%%%%%%%%%%%%%%%%%%%%%%%%%%%%%%%%%%%%
% %%%%%%%%%%%%%%%%%%%%%%%%%%%%%%%%%%%%%%%%%%%%%%%%%%%%%%%%%%%%%%%%%%%%%%%%


\begin{frame}{Distributional robustness}
\MxProfitHistogram{}

We are trying to estimate the difference in means in the MX profit data.\\[1em]

Why do we think we can do so with such a high-signal to noise ratio?\\[1em]

Only because $N = $\MxNobs{} is very large!  We can learn the effect in MX
with arbitrary accuracy --- but this does not tell us about other countries.

\end{frame}


% %%%%%%%%%%%%%%%%%%%%%%%%%%%%%%%%%%%%%%%%%%%%%%%%%%%%%%%%%%%%%%%%%%%%%%%%
% %%%%%%%%%%%%%%%%%%%%%%%%%%%%%%%%%%%%%%%%%%%%%%%%%%%%%%%%%%%%%%%%%%%%%%%%
% %%%%%%%%%%%%%%%%%%%%%%%%%%%%%%%%%%%%%%%%%%%%%%%%%%%%%%%%%%%%%%%%%%%%%%%%


\begin{frame}[t]{Distributional robustness}

Connnect AMIP with classical robustness through \emph{statistical functionals}.
%
\begin{itemize}
%
\item True, unknown distribution function $\flim(d) = p(D \le d)$
\item Empirical distribution function $\fhat(d) = \meann \ind{\d_n \le d}$
\item A statistical functional $T(F)$ (takes in $F$, returns a number).
%
\end{itemize}

Classical statistics estimates $T(\flim)$ with $T(\fhat)$.

\onslide<2->{
Sample means are an example:
%
\begin{align*}
%
T(\flim) :={}& \int \d \, \flim(\dee d)\\
\quad\quad
T(\fhat) :={}& \meann \d_n
%
\end{align*}
%
}

\onslide<3->{
Z-estimators are, too:
%
\begin{align*}
%
T(\flim) :={}& \theta\textrm{  such that  }\int G(\theta, \d) \flim(\dee d) = 0\\
T(\fhat) :={}& \theta\textrm{  such that  }\meann G(\theta, \d_n)  = 0\\
%
\end{align*}
%
}

% \only<4->{

% The influence function is just the first term in a 
% an (infinite-dimensional) Taylor series expansion.
% }

% Form an (infinite-dimensional) Taylor series expansion at some $F_0$:
% %
% \begin{align*}
% %
% T(F) = T(F_0) + T'(F_0) (F - F_0) + \mathrm{residual}.
% %
% \end{align*}
% %
% When the derivative operator takes the form of an integral
% %
% \begin{align*}
% %
% T'(F_0) \Delta = \int \infl(\x; F_0) \Delta(\dee d)
% %
% \end{align*}
% %
% then $\infl(\x; F_0)$ is known as the {\em influence function}.

% {
% Where to form the expansion? There are at least two reasonable choices:
% %
% \begin{itemize}
% %
% \item The limiting influence function $\infl(\x, \flim)$
% \item The empirical influence function $\infl(\x, \fhat)$
% %
% \end{itemize}
% %
% }
% }

    
\end{frame}




% %%%%%%%%%%%%%%%%%%%%%%%%%%%%%%%%%%%%%%%%%%%%%%%%%%%%%%%%%%%%%%%%%%%%%%%%
% %%%%%%%%%%%%%%%%%%%%%%%%%%%%%%%%%%%%%%%%%%%%%%%%%%%%%%%%%%%%%%%%%%%%%%%%
% %%%%%%%%%%%%%%%%%%%%%%%%%%%%%%%%%%%%%%%%%%%%%%%%%%%%%%%%%%%%%%%%%%%%%%%%


\begin{frame}[t]{Distributional robustness}

\def\glim{G_\infty}

Let $T(F)$ denote our estimator of the causal effect of microcredit.\\[1em]

Suppose Mexico's population is $\flim$ and our experiment gave us $\fhat$.\\
But suppose we are interested in Bolivia, with population $\flim + \Delta$!
%
\begin{align*}
\textrm{Error} ={}& T(\fhat) - T(\flim + \Delta)\\
={}&
\underbrace{T(\fhat) - T(\flim)}_{\textrm{Sampling variability }\rightarrow 0} + 
\underbrace{T(\flim) - T(\flim + \Delta)}_{\textrm{Distribution shift } \ne 0} 
\end{align*}

\textbf{How big can $\Delta$ be?}\\[1em]

\only<1>{
Classical robustness mostly considers
distribution space neighborhoods (contamination neighborhoods,
total variation distnace, KL divergence, Kolmogorov distance,
Levy-Prohorov distance, etc.) \parencite{huber1981robust,hansen2008robustness}.\\[1em]

\textbf{Can be hard for practitioners to think and talk about! }\\[1em] 

(What is the Levy-Prohorov distance between MX and Bolivia?)
}
\end{frame}


% %%%%%%%%%%%%%%%%%%%%%%%%%%%%%%%%%%%%%%%%%%%%%%%%%%%%%%%%%%%%%%%%%%%%%%%%
% %%%%%%%%%%%%%%%%%%%%%%%%%%%%%%%%%%%%%%%%%%%%%%%%%%%%%%%%%%%%%%%%%%%%%%%%
% %%%%%%%%%%%%%%%%%%%%%%%%%%%%%%%%%%%%%%%%%%%%%%%%%%%%%%%%%%%%%%%%%%%%%%%%

\begin{frame}[t]{Distributional robustness}

\textbf{How big can $\Delta$ be?}
Data dropping is simple notion of ``size'' that is
\begin{itemize}
    \item Easy to think about and 
    \item Cognizant of the task at hand.
\end{itemize}
%

\begin{minipage}{0.47\textwidth}
Define the ``perturbation--inducing proportion,''
$$
\alpha^*(\delta) = \inf\{\alpha: \textrm{AMIP}(\alpha) \ge \delta\},
$$
with $\alpha^*(\delta) = 0$ for $\delta = 0$ and $\infty$ if
the set is empty.
\end{minipage}
\begin{minipage}{0.5\textwidth}
\SimRefitOne{}
\end{minipage}

Noting that $T(\flim + \Delta) - T(\flim) \approx \int \Delta(\d) \infl(\d) \flim(\dee \d)$,
define
$$
\vnorm{\Delta} := \alpha^*\left(  \int \Delta(\d) \infl(\d) \flim(\dee \d)\right).
$$
\textbf{Intepretation:} $\vnorm{\Delta}$ = Proportion of points you need to drop to approximately
produce a change in $T$ as large as $\Delta$ does. 

\end{frame}

